\MFQTransition{事件时间仿真提示}
\begin{enumerate}[leftmargin=*]
  \item 两种策略路径的协方差越高，差值方差通常越低。
  \item 点差收入与持仓面对下一次价格变化的损益分开写。
  \item 路径为重采样单位，不能把单一路径内事件当 IID 样本。
  \item 用累计强度把事件时刻映射到近似单位指数间隔。
  \item 在改变状态前验证顺序、唯一编号和最优买卖价。
  \item 补库存、成交、费用、冲击、停止和路径分布。
  \item 成交只说明别人发生了交易，不显示你的未成交委托位置。
  \item 先完成确定性小簿，再增加随机性和真实输入。
\end{enumerate}
